\chapter{The Scan Block}
\label{ch:scan-block}

\newthought{The \yamlkey{Scan} block} answers two questions: \emph{what is this run called}, and
\emph{where do its outputs go}. It is small, but it decides the whole directory layout, so it is
worth understanding clearly.

\section{The minimal form}
\label{sec:scan-minimal}

Most cards need only two keys:

\begin{yamlwin}
Scan:
  name: "MSSM_Run"
  save_dir: "&J/outputs"
\end{yamlwin}

\begin{description}[style=nextline, leftmargin=1.2em]
  \item[\yamlkey{name}] A short, unique name for this run. It becomes the folder name that holds
        all the outputs, so prefer something filesystem-friendly (letters, digits, underscores).
  \item[\yamlkey{save\_dir}] Where output folders are created. Almost always \marker{\&J/outputs}
        --- the project-local \file{outputs/} directory. The \marker{\&J/} marker keeps the card
        portable (Chapter~\ref{ch:core-concepts}).
\end{description}

\section{What the names produce}
\label{sec:scan-layout}

\Jarvis\ builds the output tree from \yamlkey{Scan.name}. With \code{name: "MSSM\_Run"} and
\code{save\_dir: "\&J/outputs"} you get:

\begin{jfiletree}
\dirtree{%
.1 \dtfold{(project root)}.
.2 \dtfold{outputs/MSSM\_Run/}.
.3 \dtfold{DATABASE/}\DTcomment{HDF5, CSV, schema}.
.3 \dtfold{SAMPLE/}\DTcomment{per-sample files}.
.2 \dtfold{logs/MSSM\_Run/}\DTcomment{run logs}.
.2 \dtfold{images/MSSM\_Run/}\DTcomment{flowchart, plots}.
}
\end{jfiletree}

\noindent The same run name appears in \file{outputs/}, \file{logs/}, and \file{images/}, so
everything for one run is easy to find. The full output contract is described in
Chapter~\ref{ch:outputs}.

\begin{jnote}
The \file{outputs/}, \file{logs/}, and \file{images/} folders are created on first use. A fresh
project does not contain them yet.
\end{jnote}

\section{Controlling per-sample storage}
\label{sec:scan-sample-directory}

By default \Jarvis\ keeps a folder of working files for each sample under \file{SAMPLE/}. For
large scans that is a lot of folders, so the optional \yamlkey{sample\_directory} key lets you
cap and archive them:

\begin{yamlwin}
Scan:
  name: "MSSM_Run"
  save_dir: "&J/outputs"
  sample_directory:
    limit: 200
    width: 6
    archive_samples: true
\end{yamlwin}

\begin{description}[style=nextline, leftmargin=1.2em]
  \item[\yamlkey{limit}] The maximum number of live sample folders to keep at once (here, 200).
        Older ones are archived past this limit.
  \item[\yamlkey{width}] How many digits to use for the numbered folder names (here, 6, giving
        \file{000001}, \file{000002}, \ldots).
  \item[\yamlkey{archive\_samples}] When \code{true}, samples beyond the limit are compressed
        into archives instead of being kept as loose folders, keeping \file{SAMPLE/} tidy.
\end{description}

\begin{jtip}
For a quick test you can omit \yamlkey{sample\_directory} entirely. Add it once a scan grows
large enough that thousands of per-sample folders would clutter the disk. Opera-only scans can
skip per-sample folders altogether---see the \yamlkey{Runtime} options in
Chapter~\ref{ch:runtime-block}.
\end{jtip}

\section{Summary}
\label{sec:scan-summary}

\yamlkey{Scan} is two ideas: a \yamlkey{name} (which becomes the run's folder) and a
\yamlkey{save\_dir} (where that folder lives). Add \yamlkey{sample\_directory} only when you need
to control how many per-sample folders are kept. Next, the \yamlkey{Sampling} block
(Chapter~\ref{ch:sampling-block}) decides how points are drawn and scored.
